% Options for packages loaded elsewhere
% Options for packages loaded elsewhere
\PassOptionsToPackage{unicode}{hyperref}
\PassOptionsToPackage{hyphens}{url}
\PassOptionsToPackage{dvipsnames,svgnames,x11names}{xcolor}
%
\documentclass[
  letterpaper,
  DIV=11,
  numbers=noendperiod]{scrreprt}
\usepackage{xcolor}
\usepackage{amsmath,amssymb}
\setcounter{secnumdepth}{5}
\usepackage{iftex}
\ifPDFTeX
  \usepackage[T1]{fontenc}
  \usepackage[utf8]{inputenc}
  \usepackage{textcomp} % provide euro and other symbols
\else % if luatex or xetex
  \usepackage{unicode-math} % this also loads fontspec
  \defaultfontfeatures{Scale=MatchLowercase}
  \defaultfontfeatures[\rmfamily]{Ligatures=TeX,Scale=1}
\fi
\usepackage{lmodern}
\ifPDFTeX\else
  % xetex/luatex font selection
\fi
% Use upquote if available, for straight quotes in verbatim environments
\IfFileExists{upquote.sty}{\usepackage{upquote}}{}
\IfFileExists{microtype.sty}{% use microtype if available
  \usepackage[]{microtype}
  \UseMicrotypeSet[protrusion]{basicmath} % disable protrusion for tt fonts
}{}
\makeatletter
\@ifundefined{KOMAClassName}{% if non-KOMA class
  \IfFileExists{parskip.sty}{%
    \usepackage{parskip}
  }{% else
    \setlength{\parindent}{0pt}
    \setlength{\parskip}{6pt plus 2pt minus 1pt}}
}{% if KOMA class
  \KOMAoptions{parskip=half}}
\makeatother
% Make \paragraph and \subparagraph free-standing
\makeatletter
\ifx\paragraph\undefined\else
  \let\oldparagraph\paragraph
  \renewcommand{\paragraph}{
    \@ifstar
      \xxxParagraphStar
      \xxxParagraphNoStar
  }
  \newcommand{\xxxParagraphStar}[1]{\oldparagraph*{#1}\mbox{}}
  \newcommand{\xxxParagraphNoStar}[1]{\oldparagraph{#1}\mbox{}}
\fi
\ifx\subparagraph\undefined\else
  \let\oldsubparagraph\subparagraph
  \renewcommand{\subparagraph}{
    \@ifstar
      \xxxSubParagraphStar
      \xxxSubParagraphNoStar
  }
  \newcommand{\xxxSubParagraphStar}[1]{\oldsubparagraph*{#1}\mbox{}}
  \newcommand{\xxxSubParagraphNoStar}[1]{\oldsubparagraph{#1}\mbox{}}
\fi
\makeatother


\usepackage{longtable,booktabs,array}
\newcounter{none} % for unnumbered tables
\usepackage{calc} % for calculating minipage widths
% Correct order of tables after \paragraph or \subparagraph
\usepackage{etoolbox}
\makeatletter
\patchcmd\longtable{\par}{\if@noskipsec\mbox{}\fi\par}{}{}
\makeatother
% Allow footnotes in longtable head/foot
\IfFileExists{footnotehyper.sty}{\usepackage{footnotehyper}}{\usepackage{footnote}}
\makesavenoteenv{longtable}
\usepackage{graphicx}
\makeatletter
\newsavebox\pandoc@box
\newcommand*\pandocbounded[1]{% scales image to fit in text height/width
  \sbox\pandoc@box{#1}%
  \Gscale@div\@tempa{\textheight}{\dimexpr\ht\pandoc@box+\dp\pandoc@box\relax}%
  \Gscale@div\@tempb{\linewidth}{\wd\pandoc@box}%
  \ifdim\@tempb\p@<\@tempa\p@\let\@tempa\@tempb\fi% select the smaller of both
  \ifdim\@tempa\p@<\p@\scalebox{\@tempa}{\usebox\pandoc@box}%
  \else\usebox{\pandoc@box}%
  \fi%
}
% Set default figure placement to htbp
\def\fps@figure{htbp}
\makeatother

\ifLuaTeX
  \usepackage{luacolor}
  \usepackage[soul]{lua-ul}
\else
  \usepackage{soul}
\fi




\setlength{\emergencystretch}{3em} % prevent overfull lines

\providecommand{\tightlist}{%
  \setlength{\itemsep}{0pt}\setlength{\parskip}{0pt}}



 


\KOMAoption{captions}{tableheading}
\makeatletter
\@ifpackageloaded{bookmark}{}{\usepackage{bookmark}}
\makeatother
\makeatletter
\@ifpackageloaded{caption}{}{\usepackage{caption}}
\AtBeginDocument{%
\ifdefined\contentsname
  \renewcommand*\contentsname{Table of contents}
\else
  \newcommand\contentsname{Table of contents}
\fi
\ifdefined\listfigurename
  \renewcommand*\listfigurename{List of Figures}
\else
  \newcommand\listfigurename{List of Figures}
\fi
\ifdefined\listtablename
  \renewcommand*\listtablename{List of Tables}
\else
  \newcommand\listtablename{List of Tables}
\fi
\ifdefined\figurename
  \renewcommand*\figurename{Figure}
\else
  \newcommand\figurename{Figure}
\fi
\ifdefined\tablename
  \renewcommand*\tablename{Table}
\else
  \newcommand\tablename{Table}
\fi
}
\@ifpackageloaded{float}{}{\usepackage{float}}
\floatstyle{ruled}
\@ifundefined{c@chapter}{\newfloat{codelisting}{h}{lop}}{\newfloat{codelisting}{h}{lop}[chapter]}
\floatname{codelisting}{Listing}
\newcommand*\listoflistings{\listof{codelisting}{List of Listings}}
\makeatother
\makeatletter
\makeatother
\makeatletter
\@ifpackageloaded{caption}{}{\usepackage{caption}}
\@ifpackageloaded{subcaption}{}{\usepackage{subcaption}}
\makeatother
\usepackage{bookmark}
\IfFileExists{xurl.sty}{\usepackage{xurl}}{} % add URL line breaks if available
\urlstyle{same}
\hypersetup{
  pdftitle={Rapid Evidence Synthesis for Knowledge Mobilisation: A Practical User Guide},
  pdfauthor={NIHR Applied Research Collaboration Greater Manchester (ARC-GM)},
  colorlinks=true,
  linkcolor={blue},
  filecolor={Maroon},
  citecolor={Blue},
  urlcolor={Blue},
  pdfcreator={LaTeX via pandoc}}


\title{Rapid Evidence Synthesis for Knowledge Mobilisation: A Practical
User Guide}
\author{NIHR Applied Research Collaboration Greater Manchester (ARC-GM)}
\date{2027-07-03}
\begin{document}
\maketitle

\renewcommand*\contentsname{Table of contents}
{
\setcounter{tocdepth}{2}
\tableofcontents
}

\bookmarksetup{startatroot}

\chapter*{Preface}\label{preface}
\addcontentsline{toc}{chapter}{Preface}

\markboth{Preface}{Preface}

The demand for timely, reliable, and actionable evidence to inform
decision‑making in health and social care has never been greater. Across
the NHS and wider care systems, organisations face growing pressures to
adopt effective innovations, improve service delivery, and respond
rapidly to population needs. Yet the traditional processes for
identifying, appraising, and synthesising research evidence can be
time‑consuming and resource‑intensive, creating a gap between the
generation of research and its practical use.

Rapid Evidence Synthesis (RES) offers a pragmatic and responsive
approach to closing this gap for facilitating rapid knowledge
mobilisation. Developed by the NIHR Applied Research Collaboration
Greater Manchester (ARC‑GM), the RES methodology is designed to mobilise
existing evidence quickly, transparently, and in a way that remains
robust enough to support real‑world decisions. RES prioritises
relevance, clarity, and usability---focusing on what decision‑makers
want to know.

This user guide has been created to support individuals and teams who
wish to apply the ARC‑GM RES approach in their own settings. It provides
a step‑by‑step guide to planning, conducting, and communicating a Rapid
Evidence Synthesis, from the initial formulation of questions through to
interpretation and reporting. Each section includes practical tips,
examples, and tools drawn from real‑world RES projects commissioned
across the health and care system.

The user guide is intended for a wide range of users, including
researchers, knowledge mobilisation specialists, analysts, clinicians,
commissioners, and policy teams. Whether you are producing a RES
independently or working in partnership with decision‑makers, or setting
up your own RES team for producing commissioned RES, this guide aims to
enhance your confidence and capability in generating timely, meaningful,
and context‑sensitive summaries of research evidence.

Ultimately, the purpose of RES is not simply to synthesise and mobilise
knowledge, but to support informed, transparent, and evidence‑driven
decisions that can improve outcomes for patients, communities, and
health and care services. We hope that this manual equips you with the
skills and understanding needed to contribute to that goal.

To learn more about the Rapid Evidence Synthsis project of the NIHR
Applied Research Collaboration Greater Manchester (ARC‑GM), visit
\href{https://arc-gm.nihr.ac.uk/rapid-evidence-synthesis}{ARC GM
\textbar{} Rapid Evidence Synthesis}.

\bookmarksetup{startatroot}

\chapter{Overview of the Guide}\label{overview-of-the-guide}

Lord Darzi's 2024 review of the NHS in England outlines progressive
decline in NHS services {[}1{]}. In line with the review, the Government
has set out a 10 Year Health Plan for England {[}2{]}. The Plan
emphasises the urgent need for reform in response to the rising demand
for care, noting opportunities to drive improvements through the timely
adoption and spread of changes in health and social care including
innovations in digital technology {[}2{]}. It stresses that embracing
innovation is key to transforming services and delivering better
outcomes for the public. The impact of innovation relies, however, on
the use of effective approaches that offer value for money. As
recommended in the Innovation Ecosystem Programme report {[}3{]}, robust
evaluation is needed before innovations are implemented (or
de-implemented). English health and care commissioning groups such as
NHS Integrated Care Boards require timely and rapid evidence to support
change, including innovation adoption decisions, specifically:

·~~~~~~ \emph{Has this innovation, or planned change, been shown to make
a positive difference to care and/or outcomes?}

·~~~~~~ \emph{How certain is the research evidence?}

These decisions usually have cost implications, and it is essential that
NHS resources are used carefully to ensure we are selecting the right
innovations for adoption and spread.

\emph{``The capacity to generate scientific evidence {[}\ldots{}
\ldots{]} has greatly exceeded the ability to apply it to public policy
in anything but the most general way.``}

Unless actionable evidence from research is adequately mobilised, the
time and resources invested in its production are wasted and the
learning is lost to the wider health and care system. Maximising the
impact of research should be a fundamental goal.

Knowledge mobilisation (KM) enables the use of research evidence in
clinical practice and policymaking. Knowledge mobilisation brings
``\textbf{diverse communities} together to \textbf{share and create new
knowledge} in the \textbf{context of its use} to actively
\textbf{change} something'' (Knowledge Mobilisation Alliance).
Maximising research impact requires ensuring that evidence is timely,
relevant, accessible, and translated into actionable insights.

The NIHR-funded Applied Research Collaboration Greater Manchester
(ARC-GM) has developed a framework for rapid evidence synthesis (RES)
that delivers timely evidence summaries driven by the need for evidence
to support decision-making {[}4{]}. In the context of knowledge
mobilisation, the purpose of this RES approach is to summarise and
mobilise accessible research evidence for decision-makers who might
otherwise have limited capacity to access, appraise and summarise
evidence to support decision-making.

\textbf{About this RES manual and potential users}

This manual aims to guide people through the different steps involved in
producing a RES to enable them to produce own RES
\textbf{independently}, from an initial idea to the final product for
facilitating knowledge mobilisation.

The ARC-GM RES team has developed the practical manual for the following
potential users:

·~~~~~~ individuals working within health and care systems who are
interested in using RES as a knowledge mobilisation product to inform
decision-making; and/or

·~~~~~~ research or knowledge mobilisation teams within academic, NHS,
charity or other institutions who have research and/or knowledge
mobilisation capacities and are committed to producing RES to support
those delivering health and social care.

This material assumes that those considering a RES will have some prior
knowledge of how to systematically review research evidence and
knowledge of how the findings of primary research of various types can
be synthesised. Prior experience in conducting systematic reviews and
meta-analysis, however, is not required.

Those who want to explore the essentials of evidence synthesis can look
at the freely available, online resources produced by the Cochrane
Collaboration.

·~~~~~~ \href{https://www.youtube.com/watch?v=egJlW4vkb1Y}{What are
systematic reviews?}

·~~~~~~
\href{https://www.youtube.com/watch?v=WB9pbHqUs5c&list=PLgp17bk2nACMassfBMJRbfSSvpTfMgLxu&index=5}{Intro
to Systematic Reviews \& Meta-Analyses}

·~~~~~~
\href{https://www.cochrane.org/learn/courses-and-resources/evidence-essentials}{Evidence
Essentials \textbar{} Cochrane}

The manual is split into the following sections:

·~~~~~~ Section 1. Rapid Evidence Synthesis: an introduction

·~~~~~~ Section 2. Developing the RES protocol in collaboration with
decision-makers

·~~~~~~ Section 3. Searching for and selecting relevant research
evidence

·~~~~~~ Section 4. Extracting data and assessing the validity and
certainty of findings

·~~~~~~ Section 5. Synthesising and summarising the evidence

·~~~~~~ Section 6. Interpreting and communicating findings

·~~~~~~ Glossary

After reading each section you might find it helpful to complete the
tasks listed at the end to consolidate your learning.

\bookmarksetup{startatroot}

\chapter{Section 1. An Introduction to Rapid Evidence
Synthesis}\label{section-1.-an-introduction-to-rapid-evidence-synthesis}

By the end of this section you should:

·~~~~~~ Understand how the NIHR Applied Research Collaboration Greater
Manchester (ARC-GM) approach to rapid evidence synthesis (RES) compares
with other evidence synthesis approaches e.g., systematic and rapid
reviews;

·~~~~~~ Have an overview of the key steps in producing a RES;

·~~~~~~ Consider situations when generating a RES is feasible and
useful.

\textbf{Overview of evidence synthesis approaches}

In health and care systems, an \textbf{innovation} is any intentional
change in the provision of health or social care e.g.~a new treatment,
new health technology system, new test, new public health policy or
programme, or a change in working method. Health and social care
organisations make decisions to adopt and implement new innovations all
the time. Inherent in an adoption decision are considerations of
effectiveness (this change is going to make a positive difference), cost
(the cost of this change will be worthwhile) and safety (this change is
unlikely to harm service users, patients, staff or the wider system).
The kind of evidence that assists this kind of adoption decision
includes existing, relevant research evidence and it needs to be
synthesised to be given due consideration.

\emph{`Evidence synthesis refers to the process of bringing together
information from a range of sources and disciplines to inform debates
and decisions on specific issues. Decision-making and public debate are
best served if policymakers have access to the best current evidence on
an issue. An accurate, concise and unbiased synthesis of the evidence is
therefore one of the most valuable contributions the research community
can offer policymakers.'}

From `\emph{\textbf{Evidence synthesis for policy: a statement of
principles}'by the Royal Society and the Academy of Medical Sciences}
{[}5{]}.

Of the available evidence synthesis approaches, systematic review and
rapid reviews are commonly used.

`\emph{Systematic reviews seek to collate evidence that fits
pre-specified eligibility criteria in order to answer a specific
research question. They aim to minimize bias by using explicit,
systematic methods documented in advance with a protocol.}'

From \textbf{\emph{Cochrane Handbook for Systematic Reviews of
Interventions {[}6{]}.}}

See the introductory video produced by the Cochrane Collaboration
\href{https://www.youtube.com/watch?v=egJlW4vkb1Y&t=29s}{What are
systematic reviews?}

Systematic reviews are widely acknowledged to be the gold standard
approach to summarising and synthesising research findings. Reviewers
follow a pre-specified protocol (or plan), conduct an exhaustive and
comprehensive search to find research relevant to a well-defined
question, and then appraise the risk of bias in eligible studies before
synthesis (Box 1). The risk of bias assessment ensures that the results
of the research are interpreted in the light of how well the research
was done and how likely the apparent results are to be `true.'

Conducting systematic reviews is, however, often a lengthy process. A
Cochrane review takes, on average, approximately 23 months to produce
{[}7{]}. Rapid reviews emerged to address this issue, adapting
systematic review methodology `\emph{that accelerates the process of
conducting a traditional systematic review through streamlining or
omitting specific methods to produce evidence for stakeholders in a
resource-efficient manner}'{[}8{]}. So, whilst the key steps in a
systematic review are common to rapid approaches to systematic reviews
(Box 1), certain elements are curtailed. Changes include searches being
limited to a specific time period and reviewers conducting informal
critical appraisal and quality assessment. Even with these
modifications, rapid reviews can still take months to deliver.

{[}To include Box 1 here.{]}

The time it takes for systematic reviews (and rapid versions of these)
to come to fruition can prevent important research findings from
reaching decision-makers when they need them.~ Decision-makers in health
and care services often need insights from evidence in weeks rather than
months or years to inform decisions on scaling up (or scaling down)
innovations. The lack of suitable rapid evidence synthesis methods often
leads to research evidence being disregarded during decision and policy
making. This situation led us to develop our rapid approach which,
whilst imperfect, ensures that existing research evidence is part of the
process.

\textbf{The ARC-GM Rapid Evidence Synthesis (RES) process}

Our RES process generates evidence-informed insights for
decision-makers, aiming to accelerate patient and public access to
beneficial innovations. A RES can be used to help decision-making on the
adoption of all types of innovations used in the context of the health
and care system; whether this is to improve treatment, diagnosis,
education, outreach, prevention or research purposes, with the long-term
goals of improving quality, safety, outcomes, efficiency and costs.

The RES process follows the steps outlined in Box 1 but uses an approach
that can be completed in weeks rather than months or years.

RES are designed to produce a `good-enough' (rather than perfect)
summary of the evidence to support decision-making in a short timeframe,
where the alternative might be not using research at all. This approach
focuses on delivering a concise and up-to-date synthesis of the most
relevant and methodologically robust evidence available. We often
prioritise summarising pre-appraised evidence for a RES. While not
exhaustive, the RES ensures that uncertainties are clearly identified
and acknowledged.

Key considerations when developing our RES approach were speed,
transparency, an emphasis on the relevance of the evidence and analysis
of its robustness and usability.

·~~~~~~ \textbf{\emph{Speed}}. By often prioritising and summarising
existing, pre-appraised evidence synthesis in a brief report for
decision-makers, our approach can deliver a RES in 2--8 weeks, depending
on the size and scope of the motivating questions. The approach can be
applied to complex innovations and considers where evidence for
innovations may be sparse or of limited relevance.

·~~~~~~ \textbf{\emph{Transparency and robustness}}. The principles of
the \href{https://www.bmj.com/content/353/bmj.i2016}{GRADE Evidence to
Decision framework} are central to the RES approach, which makes the
RES, and the decision process to which it contributes, transparent,
consistent and reproducible.

·~~~~~~ \textbf{\emph{Relevance and usability}}. The RES approach
considers the relevance of the evidence to the local health and care
system and context as well as its reliability.

\textbf{Deciding when to conduct a RES instead of other types of
evidence synthesis}

Whilst more rapid than most other evidence synthesis options,
undertaking a RES is still resource intensive. The manual users can use
the following questions to help decide which RES requests are high
priority and feasible:

·~~~~~~ How answerable is the question for a RES? ~

·~~~~~~ How suitable is a RES and is there a better (evidence synthesis)
approach to take? When considering the review approach the risk of
getting the `wrong answer' needs to be considered. What is the tolerance
for decision-makers of drawing the wrong conclusion because the review
approach increases the risk of missing key research?

·~~~~~~ How practical and feasible is it to complete a RES? Is there an
existing evidence base for this area of inquiry?

·~~~~~~ How useful and impactful are the potential findings of a RES
likely to be in relation to decisions to adopt innovations in local
health and care systems?

Whilst a RES is often the appropriate option in the context of knowledge
mobilisation, in some cases it is not, and teams need to carefully
assess whether other evidence synthesis products like systematic reviews
are more appropriate. The decision pathway may help you determine which
type of evidence synthesis is most appropriate to undertake.

RES focuses on \textbf{\emph{knowledge mobilisation}}---that is,
summarising existing knowledge or evidence to share with specific
audiences who need it for decision-making. When the goal of an evidence
inquiry is to generate new knowledge, other types of evidence synthesis
may be more appropriate.

Here is an example of a scenario where a RES is needed as a
\emph{knowledge mobilisation} product, rather than other types of
reviews.

The Greater Manchester Integrated Care Partnership (GM ICP) Population
Health team aimed to develop regional policies to restrict the
advertising of harmful commodities in public sector--owned spaces and
wanted to use research evidence to support their practice. To make
timely and informed policy decisions, the team needed rapid access to
existing research evidence on what similar regions had already
implemented and the potential health and financial impacts of such
restrictions.

In this context, the objective was \textbf{not} to generate new
knowledge through a comprehensive and exhaustive review, but rather to
\textbf{mobilise} the most relevant, accessible, and actionable evidence
already available. Decision‑makers required a concise synthesis that
could be produced quickly, highlight what is known, identify
uncertainties, and support policy development within tight timelines. A
RES was therefore the appropriate approach, offering a pragmatic and
timely summary of the best available evidence to inform decision‑making.
For a RES, evidence may come from existing systematic reviews, primary
studies that have not yet been synthesised, and other relevant sources.
The value of producing a RES, even when systematic reviews already
exist, is that it translates the available evidence into a more
accessible format, includes critical appraisal, and assesses how
relevant the evidence is to the regional context so it can be made
actionable.

\textbf{An overview of key steps when conducting a RES}

Our RES process is intended to be flexible, allowing for a range of
questions to be addressed and using various exploratory approaches. It
takes a pragmatic approach that requires partner input for co-design and
researcher judgement to refine the approach at each stage.

The RES approach consists of several steps, which will be explained in
more detail in later sections. These steps are not necessarily linear
but iterative. The RES process is captured in the example protocol in
Box 2.

\textbf{Describe the innovation} and its purpose. Is it a public health
policy, a screening tool, a diagnostic or other test, a service, an
organisational change, or something aimed at changing individual
behaviours? What alternatives to this innovation exist? What competing
alternatives are available that could be used instead (this will include
what is being used currently)? ~

\textbf{Formulate RES questions} based on the key characteristics and
aims of the innovation. Consider the PICO(S) approach, which defines the
Population, Intervention, Comparator, Outcomes, and Settings (where
relevant) criteria that apply. Where applicable, other frameworks can be
followed (see Table \#). Initial questions often have a narrow focus to
maximise the relevance of the evidence, but we may expand our scope when
evidence is sparse (e.g., to wider innovation categories).

{[}To include the new Table here.{]}

\textbf{Consider the study design} that offers the most valid and
reliable evidence for each question (see Table 1). For questions on
innovation effectiveness, we look for systematic reviews that include
randomised controlled trials (RCTs). In the absence of existing reviews,
or if these reviews are limited or dated, we look for relevant
individual studies that have been published more recently, and may
adjust the eligibility criteria to encompass a broader range of research
evidence if necessary. ~

\textbf{Use a highly focused and tailored approach to searching} the
literature. Our literature searches are developed and conducted with an
expert where possible. We search key sources of evidence, such as the
Cochrane Library and Medline databases, and may also refer to additional
resources like the National Institute for Health and Care Excellence
(NICE) website when relevant. We stringently screen search results based
on pre-established eligibility criteria set out in the protocol (i.e.,
`rules to apply' when deciding whether to include or exclude articles).

\textbf{Extract the relevant information and evaluate evidence validity
and certainty.} Relevant information needs to be taken from the selected
material, including findings and any relevant assessment of validity and
certainty. These assessments are often guided by tools such as the
Cochrane tool for assessing the risk of bias in RCTs {[}9{]} and the
ROBINS-I tool for non-RCT evidence {[}10{]}. The GRADE approach is
widely used for assessing the certainty of a body of evidence {[}11{]}.
These judgements provide a sense of how confident we can be in the
research findings. Where existing evidence has not previously had its
validity or certainty assessed, RES reviewers may perform the relevant
assessments as needed. Additionally, at this stage we consider the
relevance of the evidence to the local context in terms of service
model, comparators and populations.

\textbf{Synthesising and summarising the evidence} to answer the key
questions. We present a clear summary of findings, including a judgement
of how confident we can be in them. We aim to explain the nature of any
uncertainty in the evidence, and how relevant it is to the key
questions.

Beyond the synthesis of evidence per se, we may also prepare the
groundwork for further evaluation (e.g., questions that should be
addressed in proof-of-value or deployment at scale), co-working with
partners.

Further information about our approach is detailed in our
\href{https://systematicreviewsjournal.biomedcentral.com/articles/10.1186/s13643-022-02106-z}{journal
publication}.

{[}To include Box 2 here{]}

\textbf{Section summary}

At ARC-GM we have developed a rapid evidence synthesis (RES) approach.
The RES approach draws on evidence synthesis methods and decision
frameworks and provides rapid assessments of existing evidence and its
relevance to specific local decision problems.

We regard RES as a suitable alternative to full evidence synthesis in
certain situations. RES serves local decision-makers who need a summary
of the best available evidence to support timely decision-making on
innovation adoption (or de-implementation). It is not always
appropriate, feasible or necessary to conduct a RES.

\textbf{Tasks to complete}

☐ Decide whether a RES is appropriate and feasible for the question you
want to answer

☐ Collaborate with your partners, if applicable

\bookmarksetup{startatroot}

\chapter{Section 2. Developing the RES protocol in collaboration with
decision-makers}\label{section-2.-developing-the-res-protocol-in-collaboration-with-decision-makers}

By the end of this section you should:

·~~~~~~ Understand the key components of a RES protocol (or plan)

·~~~~~~ Understand how to formulate decision-makers' uncertainties into
structured, answerable questions

·~~~~~~ Understand the flexible and iterative approach used when
considering eligible research for a RES.

\textbf{Involving decision-makers in RES planning}

It may be relevant to involve decision‑makers in RES planning, depending
on who the manual users are.

\textbf{\emph{1.~~~~~ If you are producing a RES for your own team:}}

As you are both the evidence user and the RES producer, you may not
necessarily need to engage additional decision‑makers.

\textbf{\emph{2.~~~~~ If you are commissioned to produce a RES by
decision‑makers:}}

In this case, decision‑makers should be actively involved in RES
planning to ensure the RES reflects their needs, priorities, and
context. Decision-makers, as primary users of the evidence, play a vital
role throughout the RES process. To ensure effective engagement and
collaboration with the partners requesting a RES, you can ask them to
fill in a pre-defined, evidence inquiry form (see our
\href{https://docs.google.com/forms/d/e/1FAIpQLSdYXNkcMp96FyfGnpNq2fGGdlUPSk2Cq8nX8fAZ5y0oyCfFaA/viewform}{\textbf{Request
for Evidence to Support-Making}} form as an example), providing
essential details about their needs and priorities. After receiving the
completed form, you can meet with the decision-makers to review the
information and discuss their requirements in greater depth.

These initial discussions focus on gathering more information on the
motivating decision problem and key questions that the decision-makers
would like to answer via the RES.

From these discussions, you can:

·~~~~~~ Identify and define relevant innovations

·~~~~~~ Propose and prioritise answerable questions. ~

This partner engagement also facilitates greater understanding of the
local and national context, ensuring the relevance of the RES to the
broader system-level perspective by considering issues such as:

·~~~~~~ where the innovations of interest are located with respect to
the wider system;

·~~~~~~ whether any system-level factors are of relevance and why.

Decision-makers may also suggest and locate relevant information and
documents that are not publicly available but which may be informative
in RES planning.

Ongoing conversations with those requesting a RES are essential to shape
a plan that reflects decision-maker needs.

We also do the following when having the initial RES planning meeting
with decision-makers {[}13{]}:

·~~~~~~ Brief the decision-makers on RES methods to help facilitate
engagement

·~~~~~~ Be honest and transparent about the desired objectives of
engagement and what a RES can and cannot do

·~~~~~~ Carefully resolve conflicts between multiple decision-makers
and/or between decision-makers and RES reviewers

·~~~~~~ Share feedback at key stages

\textbf{Overview of the key components of a RES protocol}

Preparing a RES plan involves many decisions and judgements. A
pre-specified plan is important as it guides the conduct of the RES.
This plan should describe the questions and methods relevant to each RES
and support a shared understanding of scope, whilst recognising that the
process is iterative and elements of the plan may be adapted based on
the size and shape of evidence that is identified.

We do not generally register RES protocols (e.g., on the PROSPERO
register) but we do store them internally as a record for researchers
and decision-makers to refer to as needed. ARC-GM RES protocols commonly
include the elements as shown in the example in Figure 1.

{[}\textbf{Figure 1. The ARC-GM RES plan for the topic ``policies of
restricting the use of public advertising space for unhealthy
foods''}{]}

\textbf{Description of the innovation}: When describing the target
innovations for the RES, we consider the following factors that
influence the adoption and scale-up of the innovation:

·~~~~~~ What are the characteristics of the innovation or proposed
change?

o~~ The situation(s) or setting(s) where the innovation will be used or
introduced (e.g.~public spaces, primary care, hospital-based specialist
care)

o~~ The primary purpose of the innovation (for example, preventive,
curative, rehabilitative, long-term care)

o~~ The mode of provision (e.g.~inpatient, outpatient, day care, home
care)

·~~~~~~ Who will be exposed to the innovation?

·~~~~~~ If relevant, what is the mechanism by which change will be
achieved?

·~~~~~~ If relevant, what is the strategic, political or environmental
context underlying the innovation adoption?

\textbf{Key questions:} ~The RES process involves formulating
decision-makers' queries into structured, answerable questions.
Decision-makers may have different types of questions they want to
explore in a RES, including questions about the following:

\textbf{\emph{Innovation effectiveness}}. This type of question explores
whether an innovation `works' in \textbf{treating} and/or
\textbf{preventing} a condition e.g., \emph{What is the evidence for
clinically- and cost-effective service delivery interventions for
preventing or managing multiple long-term conditions? (RES:
\href{https://arc-gm.nihr.ac.uk/media/Resources/ARC/Evaluation/RES/Multiple\%20long-term\%20conditions/ARC\%20RES\%20Multiple\%20long\%20term\%20conditions\%20-\%20final\%20report\%20(3).pdf}{Managing
multiple long-term conditions})}

\textbf{\emph{People's experiences}}. This type of question explores the
\textbf{experience} of service users or staff e.g., \emph{What evidence
is available on people's experience of debt and financial hardship
advice and support services in or via primary care settings such as
general practice and pharmacy? (RES:
\href{https://arc-gm.nihr.ac.uk/media/Resources/ARC/Evaluation/RES/Financial\%20Hardship\%20in\%20Primary\%20Care/RES\%20Financial\%20and\%20debt\%20advice\%20in\%20primary\%20care\%20-\%20new\%20template\%20FINAL.pdf}{Financial
and debt advice in primary care})}

\textbf{\emph{Barriers to and facilitators}} \textbf{\emph{of innovation
use}}. This type of question explores what inhibits and promotes the
delivery of the innovation under evaluation, e.g., \emph{What factors
are reported as barriers and facilitators for delivery of and access to
debt and financial hardship advice and support services in people,
particularly those of disadvantaged groups? (RES:
\href{https://arc-gm.nihr.ac.uk/media/Resources/ARC/Evaluation/RES/Financial\%20Hardship\%20in\%20Primary\%20Care/RES\%20Financial\%20and\%20debt\%20advice\%20in\%20primary\%20care\%20-\%20new\%20template\%20FINAL.pdf}{Financial
and debt advice in primary care})}

\textbf{\emph{Costs}}. This type of question explores whether an
innovation provides good value for what it costs e.g., \emph{What is the
evidence for clinically- and cost-effective service delivery
interventions for preventing or managing multiple long-term conditions?
(RES:
\href{https://arc-gm.nihr.ac.uk/media/Resources/ARC/Evaluation/RES/Multiple\%20long-term\%20conditions/ARC\%20RES\%20Multiple\%20long\%20term\%20conditions\%20-\%20final\%20report\%20(3).pdf}{Managing
multiple long-term conditions})}

\textbf{\emph{Diagnosis}}. This type of question explores how well a
test diagnoses someone with acondition e.g., \emph{What is the
measurement performance of remote spirometry compared with spirometry
performed in a clinical or laboratory setting? a) For patients with
suspected or diagnosed asthma? b) For patients with suspected or
diagnosed COPD? (RES:
\href{https://arc-gm.nihr.ac.uk/rapid-evidence-synthesis}{Remote
Spirometry})}

The PICO framework can help in framing a RES question, especially when
it relates to effectiveness. Other frameworks often used in framing
questions include:~

·~~~~~~ The PIT framework used for diagnostic accuracy questions (i.e.,
Population, Index test(s), Target condition),

·~~~~~~ The PICOS framework for cost-effectiveness questions (i.e.,
Population, Intervention, Comparator, Outcomes, and Setting, country, or
jurisdiction).

\textbf{Search}. This is an outline of the evidence resources and key
concepts we use to search.

\textbf{Eligibility criteria}

As with most reviews, the research question shapes the eligibility
criteria of the RES. Our RES process (unlike a typical systematic
review) requires an adaptive approach, meaning that the eligibility
criteria may be modified in response to the evidence we find. This is
informed by an adaptive process of searching and study selection, and
means that undertaking a RES is not a linear process.

Modifying the eligibility criteria might happen if, for example,
decision-makers have questions about novel innovations. In this case
there may be little or no existing research about them for the
populations of interest. So whilst we often start with a focused
approach in a RES, considering evidence that is \textbf{directly
relevant} to a question, where we find outdated, or little or no
relevant evidence we expand the eligibility criteria to look for
\textbf{indirectly relevant evidence}. Direct evidence comes from
studies that directly compare the innovations of interest among specific
populations in the context described by the RES question. Indirect
evidence does not perfectly match the specific population, innovation,
context, or outcome of interest but may provide useful insights relating
to a wider focus as described below. Decision- makers play a key role in
clarifying the types of innovation, context, and/or outcomes that are
considered directly and indirectly relevant.

Evidence could be considered indirectly relevant under the following
scenarios:

·~~~~~~ Broader classes of innovations

Whilst directly relevant evidence focuses on the same innovation the
decision problem relates to, indirectly relevant evidence might evaluate
innovations in the same `class' of interventions, such as different
version of apps with shared features. For example, when defining the
eligibility criteria for the RES of
\href{https://arc-gm.nihr.ac.uk/media/Resources/ARC/Evaluation/RES/Implementing\%20policies\%20that\%20restrict\%20the\%20use\%20of\%20outdoor\%20spaces\%20for\%20advertising\%20harmful\%20commodities/ARC-GM\%20RES\%20outdoor\%20advertising\%20harmful\%20commodities\%20revised.pdf}{Restricting
advertising in public spaces}, we initially focused on innovations of
\emph{advertising restriction policies} but, because evidence on
policies was limited, we thenexplored wider evidence on harmful
\emph{commodities' advertising}. In the RES we explained this as below:

\emph{`We acknowledge that some studies evaluated the impacts of
advertising restriction policies on outcomes whilst others evaluated the
impacts of harmful commodities' advertising on outcomes. We considered
studies that evaluated policy restrictions to be a source of direct
evidence. Where direct evidence was unavailable, we included studies
that evaluated advertising, but considered their evidence to be
indirectly relevant to this RES.' (RES:
\href{https://arc-gm.nihr.ac.uk/media/Resources/ARC/Evaluation/RES/Implementing\%20policies\%20that\%20restrict\%20the\%20use\%20of\%20outdoor\%20spaces\%20for\%20advertising\%20harmful\%20commodities/ARC-GM\%20RES\%20outdoor\%20advertising\%20harmful\%20commodities\%20revised.pdf}{Restricting
advertising in public spaces})}

·~~~~~~ Different health and care settings

A RES can be most useful when it focuses on UK-specific evidence or
evidence from other high-income countries (we often consider both these
sources to be directly relevant). When evidence is limited, however, the
eligibility criteria used for a RES may be expanded to include evidence
from low-income countries that are considered indirectly relevant. See
the example of \emph{RES:
\href{https://arc-gm.nihr.ac.uk/media/Resources/ARC/Evaluation/RES/Multiple\%20long-term\%20conditions/ARC\%20RES\%20Multiple\%20long\%20term\%20conditions\%20-\%20final\%20report\%20(3).pdf}{Managing
multiple long-term conditions}.}

·~~~~~~ Wider populations

In the example of the RES
\href{https://arc-gm.nihr.ac.uk/media/Resources/ARC/Evaluation/RES/Youth\%20workers\%20for\%20children\%20and\%20young\%20people/UoM\%20Healthier\%20Futures\%20Rapid\%20Evidence\%20Synthesis\%20-\%20Youth\%20workers\%20for\%20children\%20and\%20young\%20people\%20-\%20January\%202025.pdf}{Youth
workers for children and young people}, we aimed to answer the question
related to `\emph{children and young people who have a chronic
(long-term) physical or mental health condition}'. However, when
evidence was limited we `\emph{widened our scope and thus eligibility
criteria to include any child and young person population}' and this was
considered indirectly relevant.

As well as PICO elements, the eligibility criteria define the study
designs that will be included in the RES. As with other types of review,
RES questions dictate the most appropriate type of study to be included
(Table 1). To be as rapid as possible RES reviewers can normally limit
eligible study designs to existing, relevant systematic reviews of
appropriate primary study designs. Systematic reviews, if available,
offer reliable summaries of the relevant research. When relevant and
available, we also consider guidelines from organisations such as NICE
that make recommendations informed by robust evidence syntheses.

If suitable systematic review evidence is not available or is outdated,
we then seek directly relevant evidence from the most suitably designed
primary research studies. For clinical effectiveness questions, RES
reviewers can consider RCTs as the most robust design. Where RCTs do not
exist, you may consider other well-designed alternative quantitative
studies (see Table 1). You can ask for experts for advice when you are
unclear what study designs are the most appropriate for inclusion. In
some cases, we may also consider the findings of rapid and scoping
reviews. By incorporating other reviews and primary research, we ensure
that the RES remains as thorough and relevant as possible. ~

To help deliver a RES within the required timeline, we sometimes need to
further limit the eligibility criteria, giving details in the RES plan
based on the following:

·~~~~~~ \textbf{Quality}: We may limit inclusion to high-quality
evidence that considers factors such as study designs, methodological
rigour, and sample size.

·~~~~~~ \textbf{Comprehensiveness}: We may limit inclusion to the most
comprehensive evidence where we identify multiple systematic reviews.

·~~~~~~ \textbf{When the research was undertaken}: We may include only
the latest evidence.

As a result of applying these criteria, we prioritise the evidence that
is likely to have the most significant impact in informing policy or
practice.

{[}\textbf{Table 1. An overview of questions that a RES could answer and
the types of appropriate study designs for consideration, either as part
of systematic reviews or individually when recent, pre-appraised
evidence synthesis is not available for a RES}{]}

\textbf{Section summary}

The protocol or plan for a RES usually briefly summarises four elements:
the innovation(s) being evaluated, the key questions to be addressed,
planned approaches to searching for evidence, aand rules/criteria to
follow when deciding what evidence is eligible for inclusion.

We formulate queries from decision-makers into structured, answerable
questions, considering the elements of Participants, Interventions,
Comparators, and Outcomes (the PICO framework). We also define the types
of evidence that can be used to answer the RES question, typically using
existing systematic reviews as the primary source of evidence.

We take an iterative approach, starting with a narrow focus to find the
most directly relevant evidence, adapting eligibility criteria to widen
the scope of what can be included in the RES where evidence is limited
or narrowing what can be included (e.g., in terms of evidence quality
where there is a substantial volume of relevant evidence to consider).

\textbf{Tasks to complete:}

☐ Formulate RES question(s), using the
\href{https://livemanchesterac-my.sharepoint.com/personal/chunhu_shi_manchester_ac_uk/Documents/Documents/Research/Projects/RES/Admin/RES\%20question\%20formulation\%20form.docx}{RES
question formulation form}

☐ Specify eligibility criteria by referring to the
\href{https://livemanchesterac-my.sharepoint.com/personal/chunhu_shi_manchester_ac_uk/Documents/Documents/Research/Projects/RES/Manual/Specification\%20of\%20eligibility\%20criteria\%20form.docx}{Specification
of eligibility criteria form}

☐ Develop a RES plan, using the
\href{https://livemanchesterac-my.sharepoint.com/personal/chunhu_shi_manchester_ac_uk/Documents/Documents/Research/Projects/RES/Admin/RES\%20planning\%20tool.docx}{RES
planning tool}\textbf{\hfill\break
}

\bookmarksetup{startatroot}

\chapter{Section 3. Searching for and selecting relevant research
evidence}\label{section-3.-searching-for-and-selecting-relevant-research-evidence}

Finding research evidence that meets the eligibility criteria is a
critical part of the RES process. This section describes a streamlined
approach to identifying evidence for inclusion in a RES and provides
experience-informed insights and practical tips.

By the end of this section, you should:

·~~~~~~ Understand how to develop a search strategy and how to search
literature

·~~~~~~ Understand the streamlined approach to study selection

·~~~~~~ Understand the flexible approach to searching evidence for a RES

It is important to note that this section does not cover the know-how
and hands-on skills and steps required to perform literature searches.
Searching the literature effectively requires a certain level of skill
and experience. Although beginners can carry out searches themselves, it
is recommended to seek guidance and support from experienced librarians
where possible. Those interested in learning systematic literature
search methods are advised to refer to other resources, for example,

·~~~~~~
\href{https://training.cochrane.org/handbook/current/chapter-04}{Chapter
4: Searching for and selecting studies}, in the \emph{Cochrane Handbook
for Systematic Reviews of Interventions}.

·~~~~~~
\href{https://www.wiley.com/en-us/customer-success/cochrane/cochrane-library-user-guide}{Cochrane
Library User Guide}

·~~~~~~
\href{https://learn.nlm.nih.gov/documentation/training-packets/T0042010P/}{PubMed®
Online Training}

·~~~~~~ \href{https://hslmcmaster.libguides.com/tutorials}{Video
Tutorials} of searching the commonly used databases

\textbf{Performing and reporting a search}

RES aims to quickly mobilise relevant evidence for decision-makers on
specific health and care innovations. With this in mind, we aim to
identify the most relevant, recent and highest quality evidence rather
than conducting exhaustive searches for all potentially related
evidence.

\textbf{Constructing a suitable search strategy}

Before performing a database search, we need to develop a suitable
\textbf{search strategy}. A search strategy consists of multiple search
terms (i.e., keywords, terms, or phrases) related to a review question
that are strategically combined -- often using Boolean operators,
truncations and wildcards. Figure 2 presents the search strategy that we
used to find research for the RES
\href{https://arc-gm.nihr.ac.uk/media/Resources/ARC/Evaluation/RES/Implementing\%20policies\%20that\%20restrict\%20the\%20use\%20of\%20outdoor\%20spaces\%20for\%20advertising\%20harmful\%20commodities/ARC-GM\%20RES\%20outdoor\%20advertising\%20harmful\%20commodities\%20revised.pdf}{Restricting
advertising in public spaces}.

{[}\textbf{Figure 2. The search strategy used for Ovid MEDLINE(R){]}}

To build a search strategy, we typically follow the steps outlined
below.

\textbf{Identify keywords}: The first and vital step is identifying
keywords that can be combined in a search strategy.

In identifying keywords, we often need to:

·~~~~~~ Consider the topic(s) of the RES question(s);

·~~~~~~ Identify the main concepts within the question(s) and list
relevant keywords;

o~~~ Check whether a concept has MeSH subject headings using databases.
Subject headings are standardised terms used for precisely searching
literature on a specific topic.

·~~~~~~ Refer to key search terms used in the search strategies of
published evidence syntheses, usually available within appendices,
supplementary materials, and/or in published review protocols.

·~~~~~~ Use databases, Google, or Artificial Intelligence tools (e.g.,
\href{https://elizagrames.github.io/litsearchr/}{litsearchR}) to find
other related keywords.

For example, the terms used from line 21 to line 33 in Figure 2 are all
keywords relating to the main concept `alcohol'. We identified these
terms through multiple approaches as noted above, including referring to
the search terms used in a relevant Cochrane Review {[}15{]}.

\textbf{Use Boolean operators to combine search terms appropriately}:
Boolean operators are words (i.e., OR, AND, and NOT) used in search
engines and databases to refine and combine search terms. They help to
improve search accuracy by controlling how search terms are `connected'
with each other in a search.

·~~~~~~ OR: This word is used to combine search terms, usually for the
same concept, and increases the yield of the search by identifying
papers where the terms appear separately or together.

·~~~~~~ AND: This word is used to combine different concepts and reduces
the yield of the search by only picking papers where the 'AND' terms all
appear (ignoring papers where only one appears).

Lines 21 to 33 in Figure 2 were combined using the Boolean operator `OR'
as they are all related to the concept of alcohol.

Line 42 used the Boolean operator `AND' to combine advertising-related
terms and those related to harmful commodities such as alcohol,
gambling, and short-term loans.

\textbf{Refine a search strategy}: This often includes:

·~~~~~~ Rephrasing search terms where relevant (e.g., by using
\href{https://hslmcmaster.libguides.com/srm/truncation}{truncation and
wildcards}.

·~~~~~~ Using pre-defined search filters where available. In this
context, a search filter is a pre-defined combination of search terms
and is often validated by experts.
\href{https://sites.google.com/a/york.ac.uk/issg-search-filters-resource/home}{ISSG
Search Filters Resource} is a collection of established search filters
on various topics.

\textbf{Consideration of appropriate resources of evidence}: We
prioritise summarising pre-appraised evidence in RES in the first
instance, progressing to primary studies when this pre-appraised
evidence is not available, limited or outdated. ~We primarily search the
most commonly used databases, including Medline (or PubMed) and the
Cochrane Library, which includes both the Cochrane Database of
Systematic Reviews and the Cochrane Central Register of Controlled
Trials.

When database searching results in insufficient evidence for a RES, we
might consider using the following resources for additional searches:

·~~~~~~ \href{https://www.epistemonikos.org/en/}{Epistemonikos}, a
database of systematic reviews in health and care fields.

·~~~~~~ Resources other than databases, including evidence submitted by
stakeholders, or the manufacturers of an intervention. We often consider
\href{https://www.nice.org.uk/guidance}{NICE guidance} as a valuable
source of information and search for relevant guidance where applicable.

·~~~~~~ Reference checking or forward citation searching of relevant
evidence syntheses and primary studies.

Instead of fully presenting the detail of search methods used (e.g.,
Boolean and proximity operators used, subject headings included, text
words searched, and limits and filters applied), our RES reports only
contain brief and essential details about the search methods such as
databases and additional resources searched, the key concepts used and
dates of the searches. However, the full details of searches are
available on request. Figure 3 presents an example of the
\textbf{Search} section of the RES
\href{https://arc-gm.nihr.ac.uk/media/Resources/ARC/Evaluation/RES/Implementing\%20policies\%20that\%20restrict\%20the\%20use\%20of\%20outdoor\%20spaces\%20for\%20advertising\%20harmful\%20commodities/ARC-GM\%20RES\%20outdoor\%20advertising\%20harmful\%20commodities\%20revised.pdf}{Restricting
advertising in public spaces}

{[}\textbf{Figure 3. Example of literature searches for the RES
``Restricting
\href{https://arc-gm.nihr.ac.uk/media/Resources/ARC/Evaluation/RES/Implementing\%20policies\%20that\%20restrict\%20the\%20use\%20of\%20outdoor\%20spaces\%20for\%20advertising\%20harmful\%20commodities/ARC-GM\%20RES\%20outdoor\%20advertising\%20harmful\%20commodities\%20revised.pdf}{advertising
in public spaces}}''{]}

\textbf{Refining the search}

When starting a new RES, we design our initial search to find evidence
directly relevant to the RES questions, focusing on the specific health
and care innovations and populations that the decision-makers are
interested in. When the focused search returns a substantial amount of
evidence that is more than we can manage for a RES, we may not perform
wider searches for further evidence and indeed may narrow the
eligibility criteria further based on the quality, comprehensiveness or
timeliness of the evidence.

In cases where there is little or no evidence, or even outdated
evidence, on the specific health and care innovations and we opt to
adapt wider eligibility criteria as noted in Section 2, we expand the
search for broader, more \textbf{\emph{indirect evidence}} that does not
perfectly match the specific population, innovation, context, or outcome
of interest, but may still provide useful insights.

\textbf{Study selection}

A single reviewer normally performs study selection, although
consultation with a second reviewer is often undertaken. We often the
online tool \href{https://www.rayyan.ai/}{Rayyan} to facilitate the
study selection process.

\textbf{Section summary}

We normally follow a traditional approach to constructing a search
strategy, using Boolean operators to combine search terms, and using
commonly used resources to search for evidence.

We often take a flexible approach to searching for evidence for a RES,
beginning with a focus on direct evidence and, where necessary,
expanding to search for evidence on wider, more indirect evidence.

\textbf{Tasks to complete:}

☐ Search electronic databases and, where applicable, other resources,
referring to the
\href{https://livemanchesterac-my.sharepoint.com/personal/chunhu_shi_manchester_ac_uk/Documents/Documents/Research/Projects/RES/Manual/Development\%20and\%20execution\%20of\%20searches.docx}{Development
and execution of searches} document

·~~~~~~ Choose databases and other resources that you consider
appropriate for searches

·~~~~~~ Build appropriate search strategies

·~~~~~~ Perform searches

·~~~~~~ Download and manage search results

☐ Screen studies for inclusion against the pre-specified eligibility
criteria

☐ Document literature search and selection

\bookmarksetup{startatroot}

\chapter{Section 4. Extracting data and assessing the validity and
certainty of
findings}\label{section-4.-extracting-data-and-assessing-the-validity-and-certainty-of-findings}

Extracting relevant information and critically appraising the research
evidence and assessing its (un)certainty is a vital part of the RES
process. This section describes our streamlined approaches to these
elements of the RES and shares our practical tips.

By the end of this section, you should:

·~~~~~~ Understand the streamlined approach to extracting information
for a RES

·~~~~~~ Understand the streamlined approach to critically appraising
research for a RES

·~~~~~~ Understand different approaches to critical appraisal for
different types of research

·~~~~~~ Understand the GRADE-based approach to assessing the certainty
of research evidence by considering study validity, relevance,
consistency, and applicability.

\textbf{Extracting relevant study information for a RES}

A structured data extraction form is typically used when completing
systematic reviews. ~Once the information and data are extracted,
reviewers summarise the data and synthesise the information using
quantitative and non-quantitative approaches.

By contrast, a RES does not necessarily follow a structured data
extraction process. Instead, we adopt a streamlined
`mapping---prioritising---extracting---summarising' approach:

\begin{itemize}
\item
\item
  \textbf{Mapping}: When a RES addresses multiple questions, we group
  the included studies according to each question.
\item
\item
  \textbf{Prioritising}: If multiple studies report on the same topic or
  RES question, we aim to summarise all the relevant research. However,
  if the volume of research is unmanageable, we prioritise the most
  directly relevant, high-quality, comprehensive, and recent research,
  particularly existing evidence synthesis.
\item
\item
  \textbf{Extracting}: Following mapping and prioritisation, we extract
  key information from the included research. We focus on extracting the
  information that tells decision-makers what we have included (in terms
  of the study characteristics) and what the research evidence tells us:

  \begin{itemize}
  \item
  \item
    Author(s) and year of publication
  \item
  \item
    Basic characteristics (e.g., countries, conditions, or other
    relevant factors)
  \item
  \item
    Study design and critical appraisal
  \item
  \item
    Innovations evaluated
  \item
  \item
    Outcomes and results
  \item
  \end{itemize}
\item
\end{itemize}

Following a streamlined approach, we incorporate the key information
directly into the synthesis, presenting it in narrative form rather than
in data extraction tables.

\begin{itemize}
\item
\item
  \textbf{Summarising}: See Section 5 for more details of evidence
  synthesis approaches. Because we typically have multiple studies or
  reviews for each RES question, we often include multiple studies for
  each question, making it impractical or not helpful to present a
  narrative for the above information for every study. We therefore
  focus on the most comprehensive, high-quality and recent research,
  starting with existing systematic reviews. We then assess whether
  additional studies offer new insights or present conflicting findings.
  We draw overarching conclusions across all included studies for each
  question, using this approach.
\item
\end{itemize}

\textbf{Critical appraisal: how-to steps}

First, we examine whether the included studies have associated, existing
critical appraisal. When we include a systematic review, we often use
the critical appraisal of the included studies from the systematic
review. When no appraisal result is reported, we will perform our own
critical appraisal.

\textbf{Step 1. Understanding what needs to be appraised}

In evidence synthesis, the evidence from individual studies and the body
of the evidence produced from synthesising these individual studies are
related but distinct. Thus, we normally focus on two general issues when
deciding what needs to be appraised in a RES:

·~~~~~~ The validity of the research findings from individual studies,
i.e.~whether the design and methods used in the included research are
robust and the research findings are trustworthy.

·~~~~~~ Our confidence in the body of the total evidence reported in the
RES and our interpretation of its relevance to the motivating context
and questions.

We often begin by appraising the validity of the included studies,
followed by assessing our confidence in the evidence synthesised from
the studies.

\textbf{Step 2. Choose the appropriate critical appraisal tools}

\textbf{\emph{Assessing the validity of research results}}

In appraising the validity of the research results, we consider how
robust the included research is, assessing the risk of bias using the
appropriate tools (see Table 2). This assessment considers whether the
included research was conducted in a manner that minimises the risk of
biased results and ensures the trustworthiness of the results.

{[}\textbf{Table 2: Risk of Bias tools for different types of
evidence}{]}

When we include existing systematic reviews in a RES, we use the ROBIS
tool to appraise them. Alongside the four ROBIS domains, we may also
consider how old a review is (i.e., the date its searches were
performed) to assess whether the evidence is sufficiently contemporary.
When including primary research in a RES, we comment on the
methodological limitations following the guidance in the ROBINS-I tool
{[}10{]}. See the example in Box 3.

{[}Box 3. Example of commenting on methodological limitations in the RES
\href{https://arc-gm.nihr.ac.uk/media/Resources/ARC/Evaluation/RES/Implementing\%20policies\%20that\%20restrict\%20the\%20use\%20of\%20outdoor\%20spaces\%20for\%20advertising\%20harmful\%20commodities/ARC-GM\%20RES\%20outdoor\%20advertising\%20harmful\%20commodities\%20revised.pdf}{Restricting
advertising in public spaces}{]}

\textbf{\emph{Assessing our confidence in the evidence presented}}

This process considers the risk of bias and other factors to draw wider
conclusions about our certainty in the overall, synthesised results
being reported. This assessment follows the GRADE approach, which was
originally designed to support certainty assessment in systematic
reviews and Guidelines. More details of how to conduct GRADE assessment
are available in the
\href{https://gdt.gradepro.org/app/handbook/handbook.html}{GRADE
handbook}.

In short, the GRADE approach rates how certain we are about our
interpretation of the synthesised evidence, and the certainty is rated
using the following categories:

·~~~~~~ \textbf{high} certainty, indicating that we are confident that
the research findings reflect a true finding for that question;

·~~~~~~ \textbf{moderate} certainty, indicating that we are fairly
confident that the finding reflects a true effect;

·~~~~~~ \textbf{low} certainty, indicating that we have limited
confidence in the findings, and more research is likely to change our
conclusions;

·~~~~~~ \textbf{uncertain}, indicating that there are no clear
conclusions, based on the evidence for that question.

In the RES approach, wherever possible and appropriate, we use the GRADE
assessments already reported in the included systematic reviews. Where
we need to undertake a new GRADE assessment of our RES findings, we
employ a rapid approach to form a judgement. This assessment considers
multiple factors listed below. We apply the same GRADE criteria to
assess the certainty of evidence for each factor. Using a rapid
approach, we highlight the key issues contributing to uncertainty
without formally tabulating the assessment results of all factors. Also,
unlike the formal GRADE process, we do not consider the impact of
publication bias in our synthesis due to the rapid approaches to
searching employed.

·~~~~~~ \textbf{\emph{Study limitation}}. The methodological limitations
identified for the included studies,

·~~~~~~ \textbf{\emph{Consistency}}. Whether the different studies have
similar findings,

·~~~~~~ \textbf{\emph{Precision}}. How clear the reported effect size
is, and/or how much data (e.g., sample sizes) is used to produce the
effect size, offers a standardized measure of the magnitude of the
effect, allowing for a more meaningful interpretation of the findings.
Here the effect size is a measure people use to convey the information
about the magnitude or practical importance of the effect of a
treatment. It contains the practical information whether the effect is
large enough to be meaningful in the real world.

·~~~~~~ \textbf{\emph{Directness}}. How relevant the findings are to our
question.

See the example in the Box 4.

{[} Box 4. Example of commentary GRADE assessment in the RES
\href{https://arc-gm.nihr.ac.uk/media/Resources/ARC/Evaluation/RES/Implementing\%20policies\%20that\%20restrict\%20the\%20use\%20of\%20outdoor\%20spaces\%20for\%20advertising\%20harmful\%20commodities/ARC-GM\%20RES\%20outdoor\%20advertising\%20harmful\%20commodities\%20revised.pdf}{Restricting
advertising in public spaces} {]}

\emph{Further consideration of evidence relevance in a RES}

We usually summarise our assessment of the evidence as being of direct
or indirect relevance to the specific innovations and populations of
interest, and local or regional health and care systems.

In practice, we often consider the NHS context but may also consider
more locally relevant factors, such as whether an assessment was
conducted in a highly urbanised or rural area, and whether the study
population reflects that of a specific region as well as the UK more
broadly.

\textbf{Step 3. Document critical appraisal methods and results}

In conducting a systematic review, it is common practice to use a
structured critical appraisal form to guide and document the appraisal
process. Systematic reviews typically require detailed documentation of
the results of the critical appraisal process, including the specific
considerations used to assess each source of bias and the presentation
of assessment results in tables and/or figures.

However, rather than documenting the critical appraisal methods and
results in a detailed and structured way, our RES approach briefly
mentions the appraisal methods in the Methods section of the RES
protocol. Key study limitations identified during the appraisal are
summarised in the RES (see Section 5 for further discussion), but we do
not consider it necessary to retrieve and capture extensive information
to justify the risk of bias assessments.

\textbf{Section summary}

We often take a streamlined approach to critical appraisal and assessing
the certainty of the evidence included for a RES. This assessment
considers the validity of the included studies and the relevance of the
research results to each RES.

\textbf{Tasks to complete:}

☐ Choose critical appraisal tools that match the designs of the included
studies

☐ Critically appraise the included studies using appropriate tools

\bookmarksetup{startatroot}

\chapter{Section 5: Synthesising and summarising the
evidence}\label{section-5-synthesising-and-summarising-the-evidence}

In this section, we describe how to synthesise the evidence for a RES.

By the end of this section, you should:

·~~~~~~ Understand the narrative approach to synthesising evidence in a
RES.

·~~~~~~ Understand how available evidence and assessments of validity
and certainty are reported.

\emph{Presenting key findings of the RES}

We normally use narrative synthesis to report the evidence, mapped to
each RES question. Narrative synthesis involves descriptively
summarising and interpreting findings across studies. It is a common
method when other forms of synthesis (e.g., statistical synthesis) are
not possible. We rarely perform meta-analysis, but we often summarise
meta-analyses that have been reported in the included systematic
reviews. Meta-analysis is a way of combining the results of several
similar studies to get a more strong answer than any single study can
provide.

Different types of data (i.e., quantitative or qualitative) are
synthesised differently. For narrative syntheses of quantitative data,
our summaries of the evidence typically include the following
information:

·~~~~~~ The type and quantity of evidence (e.g., the number and type of
studies included and their sample sizes),

·~~~~~~ Effect sizes, including average effect sizes and the confidence
intervals, or numerical counts of studies that favour an innovation or
intervention,

·~~~~~~ The direction of effect sizes, which indicate whether results
favour one innovation over another, and the size and meaning of that
difference in practical or clinical terms.

This information is drawn together as in the example in Figure 5 (a
synthesis of quantitative data). In this example, information from
direct and indirect evidence is presented separately. In some cases,
this distinction is just made clear in the summary text.

These summaries of findings can also highlight the methodological
limitations of the evidence being presented, as identified in the risk
of bias assessment, and also present the level of certainty resulting
from the GRADE assessments. Incorporating the assessment of the
certainty of evidence can help gauge how likely it is that the presented
results are close to the truth.

{[}Figure 5. An example of quantitative synthesis from the RES
\href{https://arc-gm.nihr.ac.uk/media/Resources/ARC/Evaluation/RES/Implementing\%20policies\%20that\%20restrict\%20the\%20use\%20of\%20outdoor\%20spaces\%20for\%20advertising\%20harmful\%20commodities/ARC-GM\%20RES\%20outdoor\%20advertising\%20harmful\%20commodities\%20revised.pdf}{Restricting
advertising in public spaces}{]}

\emph{Presenting a summary of the RES evidence}

Each RES presents a brief summary of the evidence presented. Where there
are \ul{\textbf{multiple RES questions}}, we present the evidence
summaries across all questions more concisely (Box 5). This is like the
`abstract' of a systematic review?

{[} Box 5. Example of evidence summary in the RES
\href{https://arc-gm.nihr.ac.uk/media/Resources/ARC/Evaluation/RES/Implementing\%20policies\%20that\%20restrict\%20the\%20use\%20of\%20outdoor\%20spaces\%20for\%20advertising\%20harmful\%20commodities/ARC-GM\%20RES\%20outdoor\%20advertising\%20harmful\%20commodities\%20revised.pdf}{Restricting
advertising in public spaces} {]}

\textbf{Section summary}

In RES reports, we often use a narrative approach to synthesise
evidence, highlighting the type and quantity of evidence, the size of
effect or association, and their clinical or health significance.

Instead of extracting data using a structured form, we map the included
studies to each RES question. When several studies are available for one
question, we prioritise those that are most directly relevant,
high-quality, comprehensive, and recent. Key information is then
extracted and integrated directly into the narrative synthesis.

We summarise the evidence across all included studies and incorporate
the certainty of evidence into the interpretation of RES findings,
following a streamlined GRADE approach.

\textbf{Tasks to complete:}

☐ Map and prioritise the included studies to the RES questions; extract
sufficient information and summarise the evidence from the included
studies.

☐ Present the summary of evidence alongside the certainty of the
evidence

\bookmarksetup{startatroot}

\chapter{Section 6. Presenting and communicating
findings}\label{section-6.-presenting-and-communicating-findings}

In this section, we consider ways of presenting information in the RES
to maximise understanding of the findings and their implications.

By the end of this section, you should:

·~~~~~~ Understand how to write RES reports using clear and accessible
language

The target audience for a RES is normally policy- or decision-makers who
are frequently not research or evidence synthesis experts. The
readability of a RES is important. We outline key factors that we
believe will make a RES accessible and user-friendly.

\textbf{Clear summaries with user-friendly language and formatting:} We
always provide a summary of the evidence for each key question(s), even
when little or no (useful) evidence is identified (Box 5). Each question
is addressed separately. Where possible, we summarise existing evidence
as a narrative synthesis as detailed in Section 5 with an assessment of
certainty and relevance.

~We try to present the findings of a RES:

\begin{itemize}
\item
\item
  Using simple language, free of jargon and technical language wherever
  possible;
\item
\end{itemize}

·~~~~~~ Where using technical terms cannot be avoided, provide
explanations in simple language;

\begin{itemize}
\item
\item
  Using easy-to-understand statistical information;
\item
\item
  Using visual elements to complement text, such as tables and figures;
\item
\end{itemize}

·~~~~~~ Using an easy-to-follow layout with subheadings, bullet points,
and a consistent style.

\textbf{Concise reporting}: We usually begin a RES with a 1 to 2 page
summary of the main findings, in order to communicate the questions and
answers upfront. This summary is followed by the full report which is
typically 10 to 20 pages in length. See the examples in Box 2. We refer
you to the very useful Cochrane
\href{https://training.cochrane.org/system/files/uploads/protected_file/GUIDAN~1.PDF\#page=8}{guidance}
on writing a plain language summary along with their
\href{https://view.officeapps.live.com/op/view.aspx?src=https\%3A\%2F\%2Fwww.cochrane.org\%2Fauthors\%2Fhandbooks-and-manuals\%2Fhandbook\%2Fcurrent\%2Ftemplate-writing-cochrane-plain-language-summary.docx&wdOrigin=BROWSELINK}{template}.

\textbf{Objective tone}: A RES presents evidence that helps
decision-makers make informed choices about adopting and implementing
innovations. The report should be written in an objective tone, focusing
solely on the evidence itself. We refrain from making recommendations
about whether to adopt or reject an innovation or intervention because
we believe that research evidence is only one piece of information that
needs to be considered in decision-making.

\textbf{Section summary}

A RES report often begins with a 1--2 page summary, followed by a full
report of 10--20 pages.

We present RES findings using clear, plain language and try to keep
reports concise.

\textbf{Tasks to complete}

☐ Write up a RES report, referring to
\href{https://arc-gm.nihr.ac.uk/rapid-evidence-synthesis}{examples} of
RES reports

☐ Write a Summary section in a RES report, referring to the Cochrane
plain language summary
\href{https://view.officeapps.live.com/op/view.aspx?src=https\%3A\%2F\%2Fwww.cochrane.org\%2Fauthors\%2Fhandbooks-and-manuals\%2Fhandbook\%2Fcurrent\%2Ftemplate-writing-cochrane-plain-language-summary.docx&wdOrigin=BROWSELINK}{template}
and
\href{https://training.cochrane.org/system/files/uploads/protected_file/GUIDAN~1.PDF\#page=8}{guidance}

☐ Check the readability of the RES report

\bookmarksetup{startatroot}

\chapter*{References}\label{references}
\addcontentsline{toc}{chapter}{References}

\markboth{References}{References}

\protect\phantomsection\label{refs}




\end{document}
